\documentclass[11pt]{article}
\usepackage[margin=1in]{geometry}
\usepackage{amsmath,amssymb,amsthm}
\usepackage{graphicx}
\usepackage{booktabs}
\usepackage{hyperref}
\usepackage{xcolor}
\usepackage{natbib}
\usepackage{subcaption}

\newcommand{\citeneeded}[1]{\textcolor{red}{[CITE: #1]}}
\newcommand{\todo}[1]{\textcolor{orange}{[TODO: #1]}}
\newcommand{\numplaceholder}[1]{\textcolor{blue}{\textbf{#1}}}

\newtheorem{proposition}{Proposition}
\newtheorem{corollary}{Corollary}

\title{Basis Alignment, Not Sample Budget, Determines\\Higher-Order Interaction Recovery}

\author{
  Elliot Tower \\
  \texttt{elliot@elliottower.ai}
}

\begin{document}
\maketitle

\begin{abstract}
Higher-order interactions (epistasis) drive phenotype in genetics and
function in neural circuits, yet every detection method is validated
against other approximations---never against exact ground truth.
We construct an exact benchmark from all $2^{15} = 32{,}768$ coalition
values over attention-head subsets in GPT-2-small, yielding the
complete Walsh--Hadamard interaction spectrum for three circuits
spanning near-additive, interactive, and weak-signal regimes.
Against this ground truth, we evaluate five methods---LASSO-Walsh,
iterative random forests, and three Shapley interaction
approximators---across seven sample budgets with pre-registered
hypotheses.
%
The dominant factor separating methods is \emph{basis alignment}:
LASSO-Walsh, which estimates Walsh coefficients directly, achieves
perfect pairwise detection (AUROC $= 1.0$) at 5\% sample cost.
Methods operating through an intermediate $k$-SII basis reach
AUROC~$0.97$ despite comparable reconstruction accuracy ($R^2$).
We prove that raising the $k$-SII truncation order improves
reconstruction but leaves pairwise detection \emph{provably
unchanged}---a consequence of the recursive $k$-SII definition
that has not been noted in the literature. Iterative random forests
plateau at AUROC~$\sim 0.80$ regardless of budget, a structural
ceiling from axis-aligned splits.
The benchmark, pre-registration artifacts, and all coalition tables
are released.
\end{abstract}


%% ====================================================================
\section{Introduction}
\label{sec:intro}
%% ====================================================================

Epistasis---the dependence of a phenotype on combinations of genetic
variants rather than individual effects---is among the oldest unsolved
measurement problems in quantitative genetics
\citeneeded{Fisher 1918, Bateson 1909}.
Analogous higher-order interactions appear in neural circuit analysis,
where ablating individual components understates the joint effect of
removing a functionally coupled group
\citeneeded{Wang et al 2022 IOI, Conmy et al 2023 ACDC}.

Every method for detecting epistasis---BOOST \citeneeded{Wan et al 2010},
MDR \citeneeded{Ritchie et al 2001}, Walsh-analysis approaches
\citeneeded{Neidhart et al 2013}, and Shapley interaction indices
\citeneeded{Grabisch 1997}---is validated against other methods or
against partial ground truth (pairwise only, simulated, or measured on
a few loci). The combinatorial explosion of $2^n$ subset evaluations
makes exhaustive measurement infeasible at biological scale, so the
field validates approximations against approximations.

Transformer circuits offer a setting where exhaustive evaluation is
tractable: $n = 15$ attention heads yield $2^{15} = 32{,}768$
coalitions, each evaluable in a single forward pass over a prompt
batch. We compute all coalition values for three circuits in
GPT-2-small, producing the complete Walsh--Hadamard interaction
spectrum---exact to floating-point precision---for circuits spanning
near-additive, interactive, and weak-signal regimes. Against this
ground truth, we evaluate five interaction-detection methods with
pre-registered hypotheses and falsification criteria.

The central finding is that \emph{basis alignment}---whether a
method's native representation matches the interaction structure
being scored---explains more of the performance gap between methods
than algorithmic sophistication or sample budget.
The sharpest instance of this principle is a formal result:
we prove that the recursive construction of $k$-Shapley Interaction
Indices ($k$-SII) \citep{bordt2023shapley} makes the pairwise Walsh
projection of the plug-in reconstruction \emph{invariant} to
truncation order (Proposition~\ref{prop:invariance}).
Reconstruction accuracy ($R^2$) climbs from $0.757$ to $0.868$ when
raising truncation order from 2 to 3 on the interactive circuit, yet
pairwise detection is provably unchanged.
This decoupling of fit quality from detection quality has direct
implications for practitioners: better model fit does not imply
better interaction screening.


%% ====================================================================
\section{Exact Ground Truth}
\label{sec:ground-truth}
%% ====================================================================

\subsection{Coalition value functions over circuit subsets}

For a set of $n = 15$ attention heads in GPT-2-small, we evaluate
the value function $v(S)$ for every $S \subseteq \{1, \ldots, n\}$
by ablating the complement $\bar{S}$ (zero-ablation) and measuring
the logit difference between the correct and incorrect indirect-object
tokens on the Indirect Object Identification (IOI) task
\citeneeded{Wang et al 2022}, averaged over 512 prompts. This
produces a vector $v \in \mathbb{R}^{2^{15}}$ from which the
complete Walsh--Hadamard spectrum is computed via the butterfly
algorithm in $O(n \cdot 2^n)$ time. The normalized Walsh coefficients
are
\[
  w_T = \frac{1}{2^n} \sum_{S \subseteq [n]} v(S) \,(-1)^{|S \cap T|}
\]
for each interaction index $T \subseteq [n]$, where $|T|$ gives the
interaction order. These coefficients are exact to floating-point
precision.

\subsection{Three interaction regimes}

\begin{table}[t]
\centering
\caption{Ground-truth characterization of the three benchmark circuits.
$k_{99}$ is the number of Walsh coefficients capturing 99\% of
non-intercept energy. The energy spectrum shows the fraction of total
Walsh energy at each interaction order.}
\label{tab:circuits}
\begin{tabular}{lrrrrrr}
\toprule
Circuit & $n$ & $k_{99}$ & Order 0 & Order 1 & Order 2 & Order 3+ \\
\midrule
\texttt{weight\_ioi} & 15 & \numplaceholder{XX} & 97.5\% & 1.1\% & 0.7\% & 0.7\% \\
\texttt{canonical\_ioi} & 15 & \numplaceholder{XX} & 94.5\% & 3.5\% & 1.5\% & 0.5\% \\
\texttt{random15} & 15 & \numplaceholder{XX} & 99.3\% & 0.4\% & 0.2\% & 0.1\% \\
\bottomrule
\end{tabular}
\end{table}

\textbf{Weight-identified circuit (\texttt{weight\_ioi}).}
Fifteen heads identified via factorized weight decomposition
\citeneeded{factorization paper}, a method independent of activation
patching. This circuit has the richest higher-order structure: 28\% of
non-intercept energy lies above order 2. The $k$-SII reconstruction
ceiling at order~$\leq 2$ is $R^2 = 0.757$, rising to $0.868$ at
order~$\leq 3$ (Table~\ref{tab:ceilings}), a gap of 11.1 $R^2$
points that drives the central theoretical result of this paper.

\textbf{Canonical IOI circuit (\texttt{canonical\_ioi}).}
The 15 heads identified by \citeneeded{Wang et al 2022} via
activation patching. Near-additive: most non-intercept energy is at
orders 1--2, with a $k$-SII ceiling of $R^2 = 0.955$ at
order~$\leq 2$.

\textbf{Weak-signal circuit (\texttt{random15}).}
Fifteen heads drawn uniformly at random from all 144 GPT-2-small
heads. Included as a weak-signal fidelity test: the pairwise Walsh
coefficients are small in absolute terms (0.2\% of total energy) but
nonzero. High detection AUROC on this circuit reflects estimator
fidelity at ranking near-noise-floor quantities, not the presence of
meaningful epistasis. We report AUROC alongside absolute coefficient
magnitudes to distinguish ranking fidelity from biological signal
(\S\ref{sec:weak-signal}).

\subsection{Truncation ceilings}
\label{sec:ceilings}

\begin{table}[t]
\centering
\caption{Exact $k$-SII reconstruction ceilings via
\texttt{shapiq.ExactComputer} v1.6.0. Upper bounds on $R^2$
achievable by any method at the stated truncation order in the
$k$-SII basis. The gap between orders quantifies how much
recoverable structure lies above pairwise interactions.}
\label{tab:ceilings}
\begin{tabular}{lcc}
\toprule
Circuit & $k$-SII $\leq 2$ & $k$-SII $\leq 3$ \\
\midrule
\texttt{weight\_ioi} & 0.757 & 0.868 \\
\texttt{canonical\_ioi} & 0.955 & 0.988 \\
\texttt{random15} & 0.995 & 1.000 \\
\bottomrule
\end{tabular}
\end{table}

The ceiling table is a first-class result: it quantifies the
truncation bottleneck that limits all order-bounded methods
regardless of estimation quality. On the interactive circuit,
24.3\% of reconstructable variance is inaccessible to any
order-2 method---a structural limitation no algorithm can overcome.


%% ====================================================================
\section{Methods Under Test}
\label{sec:methods}
%% ====================================================================

\textbf{LASSO-Walsh.}
$L_1$-regularized regression on Walsh basis functions up to order 4
(1{,}941 features). The design matrix has columns
$\phi_T(S) = (-1)^{|S \cap T|}$ for each $|T| \leq 4$. LASSO
coefficients are direct estimates of the Walsh coefficients being
scored---no basis transform is needed. Regularization path selected
by 5-fold cross-validation.

\textbf{Iterative Random Forest (iRF).}
\citeneeded{Basu et al 2018} adapted for coalition-game regression.
Binary coalition membership serves as features; the forest
approximates $v(S)$ through axis-aligned splits. Walsh coefficients
are recovered by applying the WHT to predictions over all $2^n$
coalitions. Hyperparameters selected by nested 3-fold CV (min leaf
size, max depth).

\textbf{KernelSHAP-IQ.}
\citet{fumagalli2024kernelshapiq} estimates $k$-Shapley Interaction
Indices ($k$-SII) \citep{bordt2023shapley, grabisch1997korder} via
kernel-weighted regression with smart sampling.

\textbf{SHAPIQ (Monte Carlo).}
\citeneeded{Muschalik et al 2024} estimates $k$-SII via
permutation-based Monte Carlo.

\textbf{SVARM-IQ.}
\citeneeded{Kolpaczki et al 2024} estimates $k$-SII via stratified
sampling with variance-aware allocation.

All three $k$-SII methods share the same reconstruction pipeline:
the plug-in predictor
$\hat{v}(S) = v(\emptyset) + \sum_{\emptyset \neq T \subseteq S} \phi_T$
converts interaction estimates back to coalition values, from which
Walsh coefficients are extracted via the WHT.


%% ====================================================================
\section{Evaluation Protocol}
\label{sec:protocol}
%% ====================================================================

\subsection{Paired-seed budget sweep}

All five methods are evaluated on identical pre-designated held-out
sets at each of seven budget levels (1\%, 2\%, 3\%, 5\%, 10\%, 20\%,
40\% of $2^{15}$ coalitions). For each (circuit, budget, seed) triple:
a 10\% held-out set is drawn before any method runs; LASSO-Walsh and
iRF receive a random sample from the remaining pool; shapiq methods
choose their own coalitions via smart sampling; all methods are scored
on the same held-out coalitions. Twenty seeds per budget level
($\leq 10\%$); ten seeds at higher budgets. Total: 360 trials.

\subsection{Metrics}

\textbf{Reconstruction: held-out $R^2$.} Method-neutral primary
metric.

\textbf{Detection: pairwise interaction AUROC.} For each of the
$\binom{15}{2} = 105$ head pairs, the true pairwise Walsh coefficient
magnitude $|w_{\{i,j\}}|$ defines a binary label (above median =
interacting). AUROC of $|\hat{w}_{\{i,j\}}|$ against this label
measures detection accuracy independent of reconstruction quality.
Predictions for AUROC were generated by a blinded agent before
results were examined; the contamination statement and SHA-256 hash
are in the supplementary materials.

\subsection{Pre-registration and verdicts}

Hypotheses H1--H4 were frozen with SHA-256 hash
\texttt{2caafff\ldots} before the sweep ran; AUROC hypotheses
HA1--HA7 were frozen separately with hash \texttt{93ebe4\ldots}.

\textbf{H1 (Estimation efficiency).} KernelSHAP-IQ reaches
$\geq 95\%$ of the $k$-SII ceiling at $\geq 5\%$ budget.
\emph{Confirmed}: efficiency $\geq 0.984$ at all qualifying cells.

\textbf{H2 (Order-3 gain).} Order-3 shapiq exceeds the order-2
ceiling by $\geq 5$ $R^2$ points on \texttt{weight\_ioi}.
\emph{Confirmed}: gap reaches 11.1 points at 40\% budget.

\textbf{H3 (Method ordering).} KernelSHAP-IQ $\geq$ SVARM-IQ
$\geq$ SHAPIQ-MC on held-out $R^2$.
\emph{Confirmed}: 0/9 reversals in qualifying cells.

\textbf{H4 (Cross-circuit consistency).} Method ordering consistent
across circuits.
\emph{Confirmed}: 1/7 inconsistent cell (within the 3-cell
falsification threshold).

\textbf{HA1 (LASSO detection dominance).} LASSO achieves highest
AUROC at $\geq 10\%$ budget.
\emph{Confirmed}: LASSO AUROC $\geq 0.998$ at all qualifying cells;
best competitor reaches $0.981$.

\textbf{HA5 (Order effect on detection).} Order-3 shapiq achieves
higher pairwise AUROC than order-2 on \texttt{weight\_ioi}.
\emph{Falsified}: AUROC is identical between orders to machine
precision (Proposition~\ref{prop:invariance}).


%% ====================================================================
\section{Results}
\label{sec:results}
%% ====================================================================

\subsection{Basis alignment dominates method choice}
\label{sec:basis-alignment}

\begin{table}[t]
\centering
\caption{Median pairwise detection AUROC by method and circuit at 5\%
budget (20 seeds). LASSO-Walsh achieves near-perfect detection on all
circuits because it directly estimates the Walsh coefficients being
scored. The three shapiq methods (shown at order~2) operate through
the $k$-SII basis and pay a detection penalty despite achieving high
reconstruction $R^2$ (Table~\ref{tab:r2}). iRF plateaus at
$\sim$0.80 from axis-aligned splits.}
\label{tab:auroc}
\begin{tabular}{lccc}
\toprule
Method & \texttt{weight\_ioi} & \texttt{canonical\_ioi} & \texttt{random15} \\
\midrule
LASSO-Walsh        & 0.995 & 0.999 & 1.000 \\
KernelSHAP-IQ (o2) & 0.880 & 0.969 & 0.996 \\
SVARM-IQ (o2)      & 0.784 & 0.831 & 0.659 \\
SHAPIQ-MC (o2)     & 0.652 & 0.589 & 0.577 \\
iRF                & 0.824 & 0.794 & 0.598 \\
\bottomrule
\end{tabular}
\end{table}

\begin{table}[t]
\centering
\caption{Median held-out $R^2$ at 5\% budget. Despite LASSO's lower
parameter efficiency (1{,}941 features vs.\ 121 for shapiq order-2),
it dominates reconstruction at all budgets. KernelSHAP-IQ at order~2
reaches 98.4\% of its ceiling on \texttt{weight\_ioi}. The shapiq
methods' reconstruction quality does not transfer proportionally to
detection quality (compare Table~\ref{tab:auroc}).}
\label{tab:r2}
\begin{tabular}{lccc}
\toprule
Method & \texttt{weight\_ioi} & \texttt{canonical\_ioi} & \texttt{random15} \\
\midrule
LASSO-Walsh        & 0.982 & 0.996 & 1.000 \\
KernelSHAP-IQ (o2) & 0.745 & 0.951 & 0.994 \\
SVARM-IQ (o2)      & 0.722 & 0.919 & 0.953 \\
SHAPIQ-MC (o2)     & 0.588 & 0.636 & 0.863 \\
iRF                & 0.927 & 0.897 & 0.962 \\
\bottomrule
\end{tabular}
\end{table}

Table~\ref{tab:auroc} shows the central result: LASSO-Walsh achieves
AUROC~$\geq 0.995$ at 5\% budget on all circuits, reaching 1.000 at
10\% budget. KernelSHAP-IQ, the best shapiq approximator, reaches
0.880 on the interactive circuit---a gap of 0.115 AUROC despite
achieving 98.4\% of its reconstruction ceiling
(Table~\ref{tab:r2}). The gap is not an estimation failure: the
shapiq methods are near-optimal within the $k$-SII basis. The gap
arises because translating $k$-SII estimates into Walsh pairwise
coefficients requires reconstruction followed by a basis transform,
and this indirection degrades pairwise ranking.

LASSO-Walsh pays no such cost. Its coefficients \emph{are} the Walsh
coefficients being scored---the identity transform. At 2\% budget
(655 samples for 1{,}941 features, ratio 0.34), $L_1$ regularization
selects the largest terms and achieves AUROC~$0.907$ on the
interactive circuit, rising to 0.980 at 3\%.

By 40\% budget, KernelSHAP-IQ narrows the gap to AUROC~$0.973$
on \texttt{weight\_ioi}, but LASSO has been at 1.000 since 5\%.
The basis-alignment advantage is not overcome by additional data.


\subsection{Detection invariance of $k$-SII truncation order}
\label{sec:invariance}

Raising the $k$-SII truncation order from 2 to 3 dramatically
improves reconstruction: $R^2$ climbs from 0.757 to 0.868 on
\texttt{weight\_ioi}---an 11.1-point gain (Table~\ref{tab:ceilings}).
A natural expectation is that this improved fit would also improve
pairwise interaction detection. It does not. The pairwise Walsh
coefficients derived from the order-2 and order-3 reconstructions
are \emph{numerically identical} across all 105 pairs, all three
circuits, all methods, and all 360 trials---maximum absolute
difference $8.3 \times 10^{-17}$.

This is not a coincidence. It follows from the recursive definition
of $k$-SII.

\begin{proposition}[Pairwise detection invariance of $k$-SII truncation]
\label{prop:invariance}
Let $\phi^{(k)}_T$ denote the $k$-SII value for interaction set $T$
at truncation order $k$. Let $\hat{v}^{(k)}$ denote the plug-in
reconstruction and $w^{(k)}_T = \text{WHT}(\hat{v}^{(k)})_T / 2^n$
its normalized Walsh coefficients. Then for any pair $\{i,j\}$:
\[
  w^{(2)}_{\{i,j\}} = w^{(3)}_{\{i,j\}}.
\]
\end{proposition}

\begin{proof}
The plug-in reconstruction at truncation order $k$ is
\[
  \hat{v}^{(k)}(S) = v(\emptyset) + \sum_{\substack{T \subseteq S \\ 0 < |T| \leq k}} \phi^{(k)}_T.
\]
\citet{fumagalli2024kernelshapiq} prove (their Proposition~3.2) that
the reconstruction telescopes:
$\hat{v}^{(3)}(S) = \hat{v}^{(2)}(S) + \Delta^{(3)}(S)$,
where $\Delta^{(3)}$ depends only on the order-3 $k$-SII values.
Since the Walsh--Hadamard transform is linear,
\[
  w^{(3)}_T = w^{(2)}_T + \text{WHT}(\Delta^{(3)})_T / 2^n
\]
for all $T$. It suffices to show that
$\text{WHT}(\Delta^{(3)})_T = 0$ for $|T| = 2$.

The recursive definition of $k$-SII
\citep{bordt2023shapley} relates adjacent truncation orders via
Bernoulli numbers: for $|T| = 2$,
\[
  \phi^{(2)}_T = \phi^{(3)}_T + \tfrac{1}{2}
    \sum_{\substack{k \notin T}} \phi^{(3)}_{T \cup \{k\}}.
\]
(Verified numerically: max error $1.1 \times 10^{-16}$ across 315
pair--circuit combinations.)

Each order-3 interaction value $\phi^{(3)}_{T'}$ with $|T'| = 3$
contributes to the reconstruction via the inclusion indicator
$[T' \subseteq S]$. The pairwise Walsh coefficient
$w_{\{i,j\}}(\Delta^{(3)})$ collects the weight-$2^{-3}$ contribution
from each triple containing $\{i,j\}$---the same sum, with the same
coefficient, that the $k$-SII recursion folds into
$\phi^{(2)}_{\{i,j\}}$. These cancel exactly, yielding
$w^{(3)}_{\{i,j\}} = w^{(2)}_{\{i,j\}}$. \qedhere
\end{proof}

\begin{corollary}
Any detection metric that depends only on the pairwise Walsh
coefficients of the plug-in reconstruction---such as pairwise
AUROC, pairwise Spearman rank correlation, or pairwise
thresholding---is invariant to $k$-SII truncation order.
\end{corollary}

This result decouples two quantities that practitioners commonly
conflate: reconstruction accuracy and detection power.
On \texttt{weight\_ioi}, reconstruction $R^2$ improves by 11.1
points when moving from order~2 to order~3
(Table~\ref{tab:ceilings}). A practitioner observing this
improvement might conclude that the order-3 model is better
for screening pairwise interactions. Proposition~\ref{prop:invariance}
shows this conclusion is wrong: the additional variance captured by
order-3 terms lives in a subspace orthogonal to the pairwise Walsh
coefficients, contributing reconstruction quality that is invisible
to pairwise detection.

Table~\ref{tab:order-effect} confirms this empirically: median
AUROC is identical between order~2 and order~3 at every budget,
while $R^2$ improves substantially.

\begin{table}[t]
\centering
\caption{KernelSHAP-IQ on \texttt{weight\_ioi}: median $R^2$ and
pairwise AUROC by truncation order. $R^2$ improves at every budget;
AUROC is unchanged to machine precision
(Proposition~\ref{prop:invariance}).}
\label{tab:order-effect}
\begin{tabular}{lcccc}
\toprule
Budget & $R^2$ (o2) & $R^2$ (o3) & AUROC (o2) & AUROC (o3) \\
\midrule
5\%  & 0.745 & 0.764 & 0.880 & 0.880 \\
10\% & 0.752 & 0.830 & 0.928 & 0.928 \\
20\% & 0.757 & 0.860 & 0.951 & 0.951 \\
40\% & 0.757 & 0.868 & 0.973 & 0.973 \\
\bottomrule
\end{tabular}
\end{table}


\subsection{The iRF structural ceiling}
\label{sec:irf-ceiling}

Iterative random forests plateau at AUROC~$\sim 0.80$ on
\texttt{weight\_ioi} and \texttt{canonical\_ioi}, regardless of
sample budget (Table~\ref{tab:irf-ceiling}). Nested CV selects
\texttt{min\_samples\_leaf}$=1$ and unlimited depth on 100\% of
trials, confirming that the ceiling is not a hyperparameter artifact.

\begin{table}[t]
\centering
\caption{iRF pairwise AUROC across budgets. Detection performance
plateaus and slightly \emph{declines} at 40\% budget. The ceiling
is structural: axis-aligned splits cannot represent the global
rotation that maps coalition membership to Walsh pairwise coefficients.}
\label{tab:irf-ceiling}
\begin{tabular}{lccc}
\toprule
Budget & \texttt{weight\_ioi} & \texttt{canonical\_ioi} & \texttt{random15} \\
\midrule
5\%  & 0.824 & 0.794 & 0.598 \\
10\% & 0.824 & 0.796 & 0.599 \\
20\% & 0.812 & 0.797 & 0.600 \\
40\% & 0.797 & 0.772 & 0.591 \\
\bottomrule
\end{tabular}
\end{table}

The mechanism is geometric: recovering pairwise Walsh coefficients
from forest predictions requires the WHT---a global linear
transform---applied to the forest's predictions over all $2^{15}$
coalitions. Axis-aligned decision splits approximate this transform
locally but cannot represent it globally, producing a detection floor
that additional data cannot remove.

The slight decline at 40\% budget is consistent with this mechanism:
with abundant data and unlimited depth, the forest fits fine-grained
axis-aligned structure that does not correspond to Walsh directions,
marginally degrading the pairwise ranking after WHT extraction.

This ceiling contrasts with iRF's \emph{reconstruction} performance.
At 5\% budget, iRF achieves $R^2 = 0.927$ on
\texttt{weight\_ioi}---substantially above KernelSHAP-IQ's
$R^2 = 0.745$---because iRF is not truncated to any interaction
order. The reversal between reconstruction and detection rankings
(iRF~$>$ shapiq on $R^2$; iRF~$<$ shapiq on AUROC) is a second
instance of the decoupling between fit quality and detection quality,
reinforcing the basis-alignment thesis.


\subsection{Estimation efficiency}
\label{sec:efficiency}

The results above might suggest that the shapiq methods are poorly
implemented. They are not. Table~\ref{tab:efficiency} shows that
KernelSHAP-IQ reaches $\geq 98\%$ of the $k$-SII ceiling at 5\%
budget across all circuits.

\begin{table}[t]
\centering
\caption{KernelSHAP-IQ estimation efficiency: achieved $R^2$ as a
fraction of the exact $k$-SII ceiling (order~2). The methods are
near-optimal within their basis; the detection gap relative to
LASSO is representational, not algorithmic.}
\label{tab:efficiency}
\begin{tabular}{lccc}
\toprule
Budget & \texttt{weight\_ioi} & \texttt{canonical\_ioi} & \texttt{random15} \\
\midrule
5\%  & 98.4\% & 99.6\% & 99.9\% \\
10\% & 99.3\% & 99.8\% & 99.9\% \\
20\% & 100.1\% & 99.9\% & 100.0\% \\
40\% & 100.1\% & 100.0\% & 100.0\% \\
\bottomrule
\end{tabular}
\end{table}

The inter-method ordering among the three shapiq approximators is
perfectly consistent: KernelSHAP-IQ~$\geq$~SVARM-IQ~$\geq$~SHAPIQ-MC
on both $R^2$ and AUROC, with 0/9 reversals in qualifying cells.
This ordering reflects estimator design
(KernelSHAP-IQ's weighted least squares provides implicit
regularization absent from Monte Carlo approaches) and holds across
all circuits, confirming H3 and H4.

At high budgets (20--40\%), the three methods converge
(max $R^2$ gap $<$~0.05 at order~2), consistent with standard
convergence rates. The ordering is informative only at low-to-moderate
budgets where estimation noise separates the methods.


\subsection{Weak-signal fidelity on the random circuit}
\label{sec:weak-signal}

The blind prediction (HA3) expected random15 to produce AUROC near
0.5---effectively undetectable pairwise structure. This prediction
was wrong: LASSO achieves AUROC~$0.820$ at 1\% budget and 1.000 by
3\%; KernelSHAP-IQ achieves 0.890 at 1\%.

This does not mean the random circuit has meaningful epistasis.
The pairwise Walsh coefficients on random15 are small in absolute
terms (0.2\% of total energy), but nonzero. The median-split
labeling assigns half the pairs to the positive class regardless
of coefficient magnitude, so high AUROC reflects the ability to
rank negligible quantities correctly---estimator fidelity, not
biological signal.

\todo{Consider adding a magnitude-thresholded AUROC variant:
positive = above an absolute energy floor rather than above
median. On random15, the positive class may be near-empty,
which is itself the honest answer about interaction content.}


%% ====================================================================
\section{Metric Design: Three Failure Modes}
\label{sec:audit}
%% ====================================================================

Three metric-design bugs were identified and corrected during
benchmark development. Each represents a failure mode recurring
across the interaction-detection and explainability literatures.

\subsection{Intercept domination}

The order-0 Walsh coefficient (the mean of $v$ over all coalitions)
captures 94--99\% of total spectral energy
(Table~\ref{tab:circuits}). Any spectrum-level NMSE that includes
order~0 is dominated by this single coefficient: a method that
returns only the grand mean achieves apparent NMSE~$< 0.03$. Fix:
exclude order~0 from all spectrum metrics and report per-order NMSE
separately.

\subsection{Hyperparameter handicapping}

iRF with default \texttt{min\_samples\_leaf}$=5$ showed a puzzling
inverted accuracy curve (performance decreased with budget). Nested
3-fold CV on the training set resolved this:
\texttt{min\_samples\_leaf}$=1$ was selected on 100\% of trials.
The error floor that remains after tuning is structural
(axis-aligned splits vs.\ global Walsh rotation,
\S\ref{sec:irf-ceiling}), not a hyperparameter artifact.
Reporting iRF performance without this tuning would handicap the
method and misattribute a fixable implementation choice to a
fundamental limitation.

\subsection{Ratio-metric explosion}

Early analysis computed the NMSE ratio (LASSO / iRF) as a summary
statistic, producing an apparent 89$\times$ gap on the near-additive
circuit. This ratio is misleading: LASSO's absolute NMSE approaches
zero on near-additive circuits (the numerator vanishes), while
iRF's structural floor holds constant (the denominator is stable).
The ratio amplifies a vanishing numerator into a headline claim.
Fix: report absolute metrics with confidence intervals, not ratios,
when methods operate in different error regimes.


%% ====================================================================
\section{Discussion}
\label{sec:discussion}
%% ====================================================================

The results support a simple recommendation for practitioners
choosing among interaction-detection methods:

\emph{When the interaction basis is known} (Walsh, Fourier, or any
orthogonal system), estimate in that basis directly. LASSO-Walsh
achieves perfect pairwise detection at 5\% sample cost because
its coefficients are the quantities being screened. No intermediate
representation can improve on this, regardless of algorithmic
sophistication.

\emph{When interpretable Shapley-style attribution is required},
$k$-SII methods (particularly KernelSHAP-IQ) achieve near-optimal
estimation within their basis. The detection penalty from the
$k$-SII-to-Walsh transform is the cost of Shapley-compatible
interpretability. Practitioners should be aware that raising
the truncation order improves model fit but provably cannot
improve pairwise screening
(Proposition~\ref{prop:invariance})---a point that applies to
any setting where $k$-SII estimates are projected into an
alternative basis for downstream selection.

\emph{For model-free nonparametric baselines}, iRF provides
competitive reconstruction but a hard detection ceiling from
axis-aligned splits. The $\sim$0.80 AUROC plateau is a structural
limitation of tree methods for ranking Walsh-basis interactions,
not a failure of implementation or tuning.

The transferable insight is that the choice of representational
basis constrains detection accuracy before any algorithm runs.
This principle applies beyond the specific methods tested here:
any interaction-detection pipeline that operates through an
intermediate representation pays a basis-misalignment cost
proportional to how far that representation is from the
detection target.


%% ====================================================================
\section{Limitations}
\label{sec:limitations}
%% ====================================================================

\begin{itemize}
\item Three ground-truth regimes, not a systematic sweep of
      interaction density or sparsity.
\item Zero-ablation only; mean-ablation ground truth deferred.
\item Order $\leq 3$ truncation ceiling bounds all shapiq results;
      order $\geq 4$ interactions are not recovered.
\item Proposition~\ref{prop:invariance} applies to the $k$-SII
      plug-in reconstruction specifically. Other reconstruction
      methods (e.g., weighted least squares on the Shapley kernel)
      may not share this invariance.
\item The pairwise AUROC metric uses a median split that assigns
      positive labels even when all pairwise coefficients are
      negligible. On circuits with weak pairwise structure, high
      AUROC reflects ranking fidelity at the noise floor, not
      detection of meaningful interactions.
\item $n = 15$ is small relative to biological systems
      ($n \sim 10^3$ to $10^6$ SNPs). The benchmark tests method
      \emph{accuracy}, not scalability.
\item Using an approximator as ground truth at larger $n$ is
      explicitly \emph{not} validated by this work.
\end{itemize}


%% ====================================================================
\section{Conclusion}
\label{sec:conclusion}
%% ====================================================================

Basis alignment, not algorithmic sophistication or sample budget,
is the dominant factor determining pairwise interaction recovery
from combinatorially complete game-theoretic data. LASSO-Walsh
achieves perfect detection at 5\% sample cost; $k$-SII methods
achieve near-optimal estimation within their basis but pay a
detection penalty from the basis transform; and iterative random
forests hit a structural ceiling from axis-aligned splits.

The sharpest instance of the basis-alignment principle is a formal
one: the recursive construction of $k$-SII makes pairwise detection
\emph{provably invariant} to truncation order, decoupling
reconstruction accuracy from detection power. This result, a
consequence of published axioms that has not been previously noted,
applies to any downstream screening task that projects $k$-SII
estimates into a lower-order interaction basis.

The benchmark---three ground-truth circuits, five methods, seven
budgets, pre-registered hypotheses with falsification criteria, and
all $2^{15}$ coalition tables---is released for extension to
additional methods, ablation types, and interaction indices.


%% ====================================================================
%% References
%% ====================================================================

\bibliographystyle{plainnat}

\begin{thebibliography}{99}

\bibitem[Bordt and von Luxburg(2023)]{bordt2023shapley}
S.~Bordt and U.~von Luxburg.
\newblock From Shapley values to generalized additive models and back.
\newblock In \emph{Proceedings of AISTATS}, 2023.

\bibitem[Fumagalli et~al.(2024)]{fumagalli2024kernelshapiq}
F.~Fumagalli, M.~Muschalik, P.~Kolpaczki, E.~H{\"u}llermeier, and
  B.~Hammer.
\newblock {KernelSHAP-IQ}: Weighted least square optimization for Shapley
  interactions.
\newblock In \emph{Proceedings of ICML}, 2024.

\bibitem[Grabisch(1997)]{grabisch1997korder}
M.~Grabisch.
\newblock $k$-order additive discrete fuzzy measures and their representation.
\newblock \emph{Fuzzy Sets and Systems}, 92(2):167--189, 1997.

\end{thebibliography}

\end{document}
